% !TEX program = xelatex
% 由 20140521141822tgmb.dot 等版转换；使用 XeLaTeX 编译。
\documentclass[UTF8,zihao=5,fontset=none]{ctexart}
\usepackage[a4paper,top=1.35cm,bottom=1.55cm,left=1.75cm,right=1.75cm]{geometry}
\usepackage{fontspec,array,tabularx,graphicx}
\usepackage[normalem]{ulem}
\usepackage[hidelinks]{hyperref}
\usepackage{xurl}
\usepackage{setspace}
\pagestyle{empty}
\setmainfont{Times New Roman}
\setCJKmainfont[BoldFont=Songti SC Bold]{Songti SC}
\setCJKfamilyfont{song}{Songti SC}
\setCJKfamilyfont{hei}{Heiti SC}
\setCJKfamilyfont{kai}[Path=/System/Library/AssetsV2/com_apple_MobileAsset_Font7/54a2ad3dac6cac875ad675d7d273dc425010a877.asset/AssetData/]{Kaiti.ttc}
\setCJKfamilyfont{fs}[Path=/System/Library/AssetsV2/com_apple_MobileAsset_Font7/1821952872c81043711aab6910052b65da8edf2c.asset/AssetData/]{STFANGSO.ttf}
\newcommand{\song}{\CJKfamily{song}}
\newcommand{\hei}{\CJKfamily{hei}}
\newcommand{\kai}{\CJKfamily{kai}}
\newcommand{\fs}{\CJKfamily{fs}}
\setlength{\parindent}{2em}
\setlength{\parskip}{0pt}
\setlength{\tabcolsep}{4pt}
\linespread{1.03}
\newcommand{\levelone}[1]{\par\vspace{4pt}{\centering\hei\zihao{-4}#1\par}\vspace{2pt}}
\newcommand{\leveltwo}[1]{\par\vspace{2pt}{\noindent\fs\zihao{5}#1\par}}
\newcommand{\levelthree}[1]{\par\noindent{\kai\zihao{5}#1}}
\newcommand{\journalhead}{%
  \noindent\begin{tabular*}{\textwidth}{@{}l@{\extracolsep{\fill}}c@{\extracolsep{\fill}}r@{}}
  \song\zihao{5}2006 年第 1 期 & \song\zihao{5}华 南 师 范 大 学 学 报（社会科学版） & \song\zihao{5}2006 年 1 月\\[-1pt]
  \song\zihao{5}No.1,2006 & \zihao{5}JOURNAL OF SOUTH CHINA NORMAL UNIVERSITY & \zihao{5}Jan.，2006\\[-2pt]
  & \zihao{5}（SOCIAL SCIENCE EDITION） &
  \end{tabular*}\par\vspace{2pt}\hrule height 1.2pt\vspace{1.2pt}\hrule height .35pt\vspace{14pt}}
\newcommand{\titleblock}{%
  \begin{center}
  {\song\zihao{2}从监测到调控——闭环睡眠干预技术的研究进展与挑战\par}
  \end{center}\vspace{7pt}}
\newcommand{\cnabstract}{%
  {\fs\zihao{-5}\noindent 摘要：介绍了论文格式和书写，作者可以按此短文的格式排版。\hfill（小五仿）\par
  \noindent 关键词：论文；修改；格式\hfill（小五仿）\par}
  \vspace{2pt}{\song\zihao{-5}\noindent 中图分类号：\hspace{4em}文献标识码：A\hspace{3em}文章编号：1000--5455(2006)01--0000--00\par}}
\begin{document}
\journalhead
\titleblock

% BEGIN REVIEW BODY
\levelone{一、引\quad 言}

\leveltwo{（一）睡眠障碍的疾病负担与睡眠干预的临床需求}

睡眠障碍是具有公共健康意义的系统性社会健康挑战。欧洲 47 国整合分析显示，成人阻塞性睡眠呼吸暂停（obstructive sleep apnea，OSA）、失眠、不宁腿综合征（restless legs syndrome，RLS）、发作性睡病和快速眼动（rapid eye movement，REM）睡眠行为障碍患病率分别为 18\%、10\%、3\%、0.03\% 和 0.009\%\textsuperscript{[1]}。失眠的诊断级负担随判定口径而变化：面对面《精神障碍诊断与统计手册》（Diagnostic and Statistical Manual of Mental Disorders，DSM）访谈下合并患病率为 12.4\%，成人临床相关慢性失眠全球估计为 16.2\%、约 8.52 亿人\textsuperscript{[2--3]}。中国 OSA 按呼吸暂停低通气指数（apnea-hypopnea index，AHI）$\geq 5$ 的汇总患病率为 11.8\%，全球 20--79 岁成人 RLS 模型化患病率为 7.12\%\textsuperscript{[4--5]}。这些结果共同说明，睡眠干预需求受疾病类别、地区范围和诊断口径共同塑造，不能被合并为单一负担数值。

在可穿戴人工智能睡眠技术领域，现有研究主要集中于筛查与诊断；一项纳入 46 项研究的范围综述没有发现将可穿戴人工智能用于治疗的研究，表明从状态监测与识别到治疗性干预的转化链条仍缺少直接研究证据\textsuperscript{[6]}。在成人慢性失眠的规范治疗方面，欧洲指南将失眠认知行为治疗（cognitive behavioral therapy for insomnia，CBT-I）列为一线治疗，美国睡眠医学会（American Academy of Sleep Medicine，AASM）指南也强推荐多组分 CBT-I\textsuperscript{[7--8]}。尽管 CBT-I 已被指南确立为循证治疗，其现实应用仍受到服务可及性有限以及部分患者退出或依从性不足的制约\textsuperscript{[8--10]}。这些实施障碍与技术转化缺口并不等同于闭环干预已具疗效：慢性失眠闭环听觉刺激试点虽观察到慢振荡即时增强，却未改善记忆巩固、宏观睡眠结构、觉醒次数或主观睡眠质量\textsuperscript{[11]}。这一结果表明，实现特定神经生理靶点的实时调控，并不意味着其能够稳定转化为具有临床意义的睡眠改善或日间获益。因此，整合连续感知、状态识别、即时干预与反馈调整的闭环系统，更适合被视为改善监测与治疗衔接的重要候选技术路径之一；其临床价值仍有赖于刺激算法个体化、长期疗效及患者重要结局的进一步验证。

\leveltwo{（二）从睡眠监测到睡眠调控：闭环范式的提出}

从“睡眠监测”向“睡眠调控”的过渡，核心在于将连续状态感知实时转化为可更新的干预决策。离线监测主要回答整夜睡眠结构与事件分布，实时识别进一步判断当前睡眠状态，而闭环调控还需依据当前脑状态决定是否实施刺激，并根据刺激后的生理响应更新后续决策。与从随机慢振荡相位启动、随后按预设间隔递送刺激的开放环方案相比，闭环系统强调干预时机与目标神经生理状态的匹配\textsuperscript{[12]}。在严格的术语边界下，只有持续监测目标状态，并利用反馈信息维持或调节该状态的方案，才可称为闭环状态依赖刺激；因此，并非所有依据睡眠阶段触发的刺激均构成完整闭环\textsuperscript{[13]}。

这一范式转变相应拓展了睡眠脑机接口的算法任务边界。算法所面向的控制对象既包括睡眠阶段，也包括慢振荡、睡眠纺锤波等瞬态事件及更精细的振荡相位信息；因此，相关算法的功能需从离线睡眠分期扩展至在线状态识别、瞬态事件检测、相位预测、刺激时机决策及闭环反馈控制。Wireless Dreem Device（WDD）将非快速眼动第 3 期（N3）门控、慢振荡相位估计和声音刺激递送串联起来，并在与多导睡眠监测（polysomnography，PSG）同步记录的健康青年样本中，实现了灵敏度为 0.70、特异度为 0.90 的实时 N3 检测；同时，该系统能够将声音刺激定位于慢振荡上升相附近，刺激相位为 $45\pm52^{\circ}$。这一结果初步表明，可穿戴脑电图（electroencephalography，EEG）设备具备在居家环境中实施状态依赖性睡眠刺激的技术可行性\textsuperscript{[14]}。

由此，闭环系统的评价标准也由离线识别性能转向端到端调控性能。Ferster 等将非快速眼动睡眠（non-rapid eye movement，NREM）、慢波活动（slow-wave activity，SWA）、beta 功率和目标相位纳入联合门控；Portiloop 则进一步将采集、滤波、推理、决策和刺激执行的累计时延以及误触发、漏触发纳入工程评价\textsuperscript{[15--16]}。因此，离线准确率不能单独代表闭环性能，还需综合考察端到端时延、触发可靠性、目标相位门控和扰动保护等实时约束。现有研究在目标事件、设备平台、研究人群和时延口径方面仍存在明显异质性，尚不足以进行直接的系统间比较。可穿戴 EEG 为长期居家闭环调控提供了技术条件，但系统可执行性和靶点调控能力并不等同于临床疗效，这也构成后续分析闭环系统架构及其临床验证路径的基本前提。

\leveltwo{（三）闭环睡眠干预系统的基本架构}

闭环睡眠干预的基本架构可理解为由感知、识别/决策、刺激执行、反馈与失败保护共同组成的在线系统；在闭环听觉刺激（closed-loop auditory stimulation，CLAS）框架中，其最小实现包括实时脑活动访问、目标振荡检测和刺激系统\textsuperscript{[13]}。感知层首先要保证输入信号可用，例如 ENMod 以 5 秒均方根（root mean square，RMS）选择额区干 EEG 通道，但其居家研究中约 30\% 数据集仍无可用 EEG，提示电极接触与质量控制会直接限制后续识别和干预环节\textsuperscript{[17]}。在识别/决策层，系统可联合 NREM、SWA、beta 功率和目标相位进行门控，并在运动、睡眠阶段变化或觉醒时拒绝或停止触发\textsuperscript{[14--15]}。执行层则需同时评估采集、滤波、推理、决策和输出延迟，以及误触发与漏触发；Portiloop 的纺锤波工程案例报告 64 ms 常数延迟\textsuperscript{[16]}。从控制系统的严格定义看，依据当前脑状态触发预设刺激主要属于事件触发式闭环；只有在刺激后继续测量生理或行为响应，并利用该反馈动态调整后续刺激的时机、强度或其他参数，才构成具有在线自适应能力的完整反馈控制。上述术语边界借鉴自非睡眠脑电图触发经颅磁刺激（electroencephalography-triggered transcranial magnetic stimulation，EEG-TMS）闭环研究。就现有睡眠研究证据而言，上述系统证据主要支持“检测—触发”意义上的闭环运行；“刺激—反馈—参数更新”这一自适应控制链条仍需在睡眠场景中得到完整实现和充分验证\textsuperscript{[18--19]}。

\leveltwo{（四）本文综述范围与结构安排}

基于上述临床需求、范式边界与系统架构，本文将闭环睡眠干预界定为一种以睡眠相关状态或事件的实时感知为前提，并将识别结果用于干预触发、参数调节或效果评估的技术体系。后文将首先讨论可穿戴传感、实时信号处理、睡眠分期、事件/相位估计以及端到端时延等关键技术环节；随后围绕慢波睡眠和全睡眠周期两类应用场景，梳理听觉刺激、电/磁刺激及多模态刺激系统的实现方式、神经生理效应、行为或功能终点与验证层级；最后总结该领域在从生理增强走向临床获益、个体化参数设定、长期安全性以及居家真实世界评测等方面面临的挑战。本文不纳入缺乏实时触发逻辑、未设置睡眠相关终点，或仅有产品宣称而缺少可核验证据支持的研究。

\vspace{3pt}\noindent{\hei 参考文献：}\par
{\song\zihao{-5}\setlength{\parindent}{0pt}\setlength{\parskip}{1pt}
[1]\quad BASSETTI C L A, WELTER L S, MONTES-MARTINEZ M, et al. Epidemiology and economic burden of sleep disorders in Europe. European Journal of Neurology, 2026, 33(2): e70463.\par
[2]\quad VAN STRATEN A, WEINREICH K J, F\'ABIAN B, et al. The prevalence of insomnia disorder in the general population: A meta-analysis. Journal of Sleep Research, 2025, 34(5): e70089.\par
[3]\quad BENJAFIELD A V, SERT KUNIYOSHI F H, MALHOTRA A, et al. Estimation of the global prevalence and burden of insomnia in adults: A systematic review and modelling study. Sleep Medicine Reviews, 2025, 82: 102121.\par
[4]\quad NIU Y, SUN S, WANG Y, et al. Spatiotemporal trends in the prevalence of obstructive sleep apnea in China from 2000 to 2024: A systematic review and meta-analysis. Nature and Science of Sleep, 2025, 17: 879--903.\par
[5]\quad SONG P, WU J, CAO J, et al. Global and regional prevalence of restless legs syndrome among adults: A systematic review and modelling analysis. Journal of Global Health, 2024, 14: 04113.\par
[6]\quad AZIZ S, ALI A, ASLAM H, et al. Wearable artificial intelligence for sleep disorders: Scoping review. Journal of Medical Internet Research, 2025, 27: e65272.\par
[7]\quad RIEMANN D, BAGLIONI C, BASSETTI C, et al. European guideline for the diagnosis and treatment of insomnia. Journal of Sleep Research, 2017, 26(6): 675--700.\par
[8]\quad EDINGER J D, ARNEDT J T, BERTISCH S M, et al. Behavioral and psychological treatments for chronic insomnia disorder in adults: An American Academy of Sleep Medicine clinical practice guideline. Journal of Clinical Sleep Medicine, 2021, 17(2): 255--262.\par
[9]\quad KOFFEL E, BRAMOWETH A D, ULMER C S. Increasing access to and utilization of cognitive behavioral therapy for insomnia (CBT-I): A narrative review. Journal of General Internal Medicine, 2018, 33(6): 955--962.\par
[10]\quad MATTHEWS E E, ARNEDT J T, MCCARTHY M S, et al. Adherence to cognitive behavioral therapy for insomnia: A systematic review. Sleep Medicine Reviews, 2013, 17(6): 453--464.\par
[11]\quad DUDYSOV\'A D, JANK\r{U} K, PIORECK\'Y M, et al. Closed-loop auditory stimulation in chronic insomnia disorder: A pilot study. Journal of Sleep Research, 2024, 33(6): e14179.\par
[12]\quad WEIGENAND A, M\"OLLE M, WERNER F, et al. Timing matters: Open-loop stimulation does not improve overnight consolidation of word pairs in humans. European Journal of Neuroscience, 2016, 44(6): 2357--2368.\par
[13]\quad ESFAHANI M J, FARBOUD S, NGO H-V V, et al. Closed-loop auditory stimulation of slow oscillations during sleep: Principles and applications. Neuroscience \& Biobehavioral Reviews, 2023, 153: 105379.\par
[14]\quad DEBELLEMANIERE E, CHAMBON S, PINAUD C, et al. Performance of an ambulatory dry-EEG device for auditory closed-loop stimulation of sleep slow oscillations in the home environment. Frontiers in Human Neuroscience, 2018, 12: 88.\par
[15]\quad FERSTER M L, DA POIAN G, MENACHERY K, et al. Benchmarking real-time algorithms for in-phase auditory stimulation of low-amplitude slow waves with wearable EEG devices during sleep. IEEE Transactions on Biomedical Engineering, 2022, 69(9): 2916--2925.\par
[16]\quad VALENCHON N, BOUTEILLER Y, JOURDE H R, et al. The Portiloop: A deep learning-based open science tool for closed-loop brain stimulation. PLOS ONE, 2022, 17(8): e0270696.\par
[17]\quad BRESSLER S, NEELY R, YOST R M, et al. A wearable EEG platform for closed-loop stimulation during sleep. Journal of Neural Engineering, 2023, 20(5): 056030.\par
[18]\quad BERGMANN T O. Brain state-dependent brain stimulation. Frontiers in Psychology, 2018, 9: 2108.\par
[19]\quad ZRENNER C, ZIEMANN U. Closed-loop brain stimulation. Biological Psychiatry, 2024, 95(6): 545--552.\par}
% END REVIEW BODY

\end{document}
